\clearpage
\fancyhead[L]{\small MOLLWEIDE. COMMENTATIONES.}
\hypertarget{native-572-lower}{}
\begin{center}\small Göttingische gelehrte Anzeigen. Stück 71. Seiten 705..710. 2. Mai 1814.\\[3pt]\rule{.72\linewidth}{.3pt}\end{center}

\emph{Commentationes mathematico-philologicae tres, sistentes explicationem duorum locorum difficilium, alterius Virgilii, alterius Platonis itemque examinationem duorum mensurarum praeceptorum Columellae. Adjecta est epistola ad v. cl.} J. G. \textsc{Schneider} \emph{de excerptis geometricis Epaphroditi et Vitruvii Rufi scripta ab auctore harum commentationum} \textsc{Carolo Brandano Mollweide}, astron. in acad. Lipsiensi professore. 1813. Leipzig. Bei Cnobloch. 122 Seiten im Oktavformat, dazu eine Kupfertafel.

Je seltener gründliche mathematische Kenntnisse und eine eingehende Vertrautheit mit den Schriftstellern und Sprachen des Altertums in einer Person zusammenkommen, desto willkommener müssen die Aufschlüsse sein, die der gelehrte Verfasser der vorliegenden Aufsätze mithilfe der Mathematik zu verschiedenen dunklen Stellen bei antiken Schriftstellern gibt. Der erste Aufsatz betrifft die bekannte Stelle in Vergils Gedicht über den Landbau zu den Zeiten der beiden Honigernten.

\SourceLeaf{573}\begin{quote}\small Bis gravidos cogunt fetus: duo tempora messis \\ Taygete simul os terris ostendit honestum \\ Pleias, et oceani spretos pede repulit amnes, \\ Aut eadem sidus fugiens ubi piscis aquosi \\ Tristior hibernas coelo descendit in undas.

\hfill Georg. IV. 231---235. \end{quote}

Die meisten Ausleger sind sich einig, dass in den letzten beiden Versen der Herbst gemeint ist, wenn die Plejaden morgens untergehen, und dass daher mit descendit in undas der \emph{kosmische} Untergang gemeint sein muss. Für die Worte sidus fugiens piscis aquosi hat man jedoch eine Fülle verschiedener Erklärungen versucht, ohne dass eine davon allgemeine Zustimmung gefunden hätte. Einige Ausleger haben offen eingeräumt, dass die Worte so, wie sie dastehen, unerklärlich seien. Herr \textsc{Mollweide} hatte seine Auffassung dieser Stelle schon vor geraumer Zeit in der Monatl.\ Corresp. (Mai 1802) vorgetragen. Später legte er anlässlich seiner Ernennung zum Professor der Astronomie in Leipzig in einer eigens gedruckten kleinen Schrift eine Umarbeitung jenes Aufsatzes vor, die jetzt mit weiteren Ergänzungen in der vorliegenden Sammlung erscheint. Unter piscis aquosus versteht er im Unterschied zu den meisten anderen Auslegern, die darin die Fische des Tierkreises sehen, den Südlichen Fisch. Allerdings hat er in den späteren Umarbeitungen einen der Hauptgründe fallen lassen, mit denen er diese Auffassung im ersten Aufsatz gestützt hatte (das Beiwort aquosus, das bei den Fischen des Tierkreises etwas überflüssig zu sein scheine, beziehe sich nämlich auf die Stellung des Sterns erster Größe im Südlichen Fisch am Wasserstrahl des Wassermanns). Vermutlich tat er das, weil es andere Stellen gibt, an denen piscis mit demselben Beiwort eindeutig den Fisch des Tierkreises bezeichnet. Jetzt versteht er dieses Beiwort nur als Anspielung auf die nasse Jahreszeit. Das \emph{Fliehen} der Plejaden vor dem Südlichen Fisch erklärt er damit, dass der kosmische Untergang der Plejaden und der akronychische Aufgang des Südlichen Fisches zeitlich ziemlich nahe beieinanderliegen. Zum Beweis, dass das Bild einer Flucht in diesem Sinne gebraucht werden kann, verweist er auf eine andere Stelle desselben Gedichts.

\enlargethispage{-10\baselineskip}
\begin{quote}\small Vere fabis satio; tum te quoque, Medica, putres \\ Accipiunt sulci, et milio venit annua cura, \\ Candidus auratis aperit quum cornibus annum \\ Taurus, et adverso \emph{cedens} canis occidit astro.

\hfill Georg. I. 215---218. \end{quote}

\noindent Hier hat cedens dieselbe Bedeutung wie fugiens an der obigen Stelle und deutet nach Herrn \textsc{Mollweide} an, wie nahe der Abenduntergang des Hundes und der morgendliche Aufgang des Stiers zeitlich beieinanderliegen. So viel Gelehrsamkeit Herr \textsc{Mollweide} auch zur Stütze seiner Erklärung aufbringt, so kann der Rezensent doch nicht leugnen, dass sie auf ihn immer noch etwas gezwungen wirkt. Er lässt diese Erklärung der zweiten Stelle gern gelten. Wenn der Abenduntergang eines Sterns mit dem morgendlichen Aufgang eines anderen zusammentrifft und dies ein Fliehen des ersten vor dem zweiten genannt wird, erkennt er darin \emph{trotz der verschiedenen} Tageszeiten ein schönes und wahres poetisches Bild. Dasselbe Bild auf den umgekehrten Fall anzuwenden (um den es hier geht), erscheint ihm jedoch hart. Hier geht der fliehende Stern morgens unter, und der Stern, der ihn vertreibt, geht abends auf. Im Frühjahr beginnt das Sternbild des Großen Hundes, das bis dahin jeden Abend am westlichen Himmel geglänzt hatte, ganz unsichtbar zu werden, etwa zu der Zeit, zu der der Stier, der seinerseits eine Weile ganz unsichtbar gewesen war, wieder sichtbar zu werden beginnt und in den frühen Morgenstunden immer mehr den östlichen Himmel schmückt. Was kann man dagegen haben, wenn ein \SourceLeaf{574}Dichter dies ein Fliehen des Hundes vor dem Stier nennt? Aber was bleibt von der Wahrheit und Anschaulichkeit des poetischen Bildes, wenn der Dichter die Plejaden in einer Jahreszeit vor Fomahand fliehen lässt, in der die Erscheinungen tatsächlich in folgender Reihenfolge eintreten? Abends nach Sonnenuntergang, sobald überhaupt Sterne sichtbar werden, sind beide schon aufgegangen, Fohamond freilich fast eine halbe Stunde früher als die Plejaden. Doch Fomahand steigt nur wenig über den Horizont und ist lange vor Mitternacht schon wieder untergegangen, während die Plejaden die ganze Nacht hoch am Himmel glänzen und erst in dem Augenblick unter den Horizont sinken, in dem die Sonne aufgeht. So scheint es, dass die obige Stelle, sofern man sie als Parallelstelle heranziehen darf, eher die Erklärung des Bischofs \textsc{Horsley} stützen würde, der in den letzten beiden Versen den Abenduntergang der Plejaden und den morgendlichen Aufgang der Fische (im Frühjahr) als gleichzeitig ansieht, wenn dieser Auslegung nicht andere, nicht unerhebliche Bedenken im Wege stünden. [Die Besprechungen der zweiten und dritten Abhandlung und die Briefe finden sich in Band IV dieser Werke unter den Anzeigen von Schriften, die Gegenstände aus dem Gebiet der Geometrie behandeln].
